\documentclass[a4paper,12pt]{article}

\usepackage[T2A]{fontenc}
\usepackage[utf8]{inputenc}
\usepackage[russian]{babel}
\usepackage{amsmath,amssymb,mathtools}
\usepackage{helvet}
\renewcommand{\familydefault}{\sfdefault}
\usepackage{icomma}
\usepackage[a4paper,margin=20mm]{geometry}
\usepackage{microtype}
\usepackage{tikz}
\usetikzlibrary{calc,arrows.meta,positioning,shapes.geometric}
\usepackage{tabularx,booktabs,longtable}
\usepackage{enumitem}
\usepackage{fancyhdr}
\usepackage{setspace}
\usepackage{xcolor}
\usepackage{hyperref}
\usepackage{parskip}

\definecolor{accent}{HTML}{2D6A6A}
\definecolor{darkaccent}{HTML}{1E4D4D}
\definecolor{textgray}{HTML}{303030}
\definecolor{softgray}{HTML}{6C757D}
\definecolor{lightaccent}{HTML}{E8F1F1}

\hypersetup{colorlinks=true,linkcolor=darkaccent,urlcolor=darkaccent}
\onehalfspacing
\setlength{\parindent}{0pt}
\setlist[itemize]{leftmargin=1.5em,itemsep=0.35em,topsep=0.4em}
\setlist[enumerate]{leftmargin=1.7em,itemsep=0.55em,topsep=0.4em}

\pagestyle{fancy}
\fancyhf{}
\fancyhead[L]{\small\color{softgray}Прикладные задачи}
\fancyhead[R]{\small\color{softgray}Матвей Горбачев}
\fancyfoot[C]{\small\thepage}
\renewcommand{\headrulewidth}{0.3pt}
\renewcommand{\footrulewidth}{0pt}
\setlength{\headheight}{15pt}

\newcommand{\sectionline}{\par\vspace{0.3em}{\color{accent}\hrule height 0.7pt}\vspace{0.7em}}
\newcommand{\key}[1]{\textbf{\color{darkaccent}#1}}

\tikzset{
  startstop/.style={ellipse,draw=accent,thick,fill=lightaccent,minimum width=38mm,minimum height=10mm,align=center,font=\small},
  process/.style={rectangle,rounded corners=2mm,draw=accent,thick,minimum width=105mm,minimum height=11mm,align=center,font=\small,text width=95mm},
  decision/.style={diamond,aspect=2.3,draw=accent,thick,align=center,font=\small,text width=38mm,inner sep=1.5mm},
  arrow/.style={-{Latex[length=3mm]},thick,draw=darkaccent}
}

\begin{document}

\begin{titlepage}
\thispagestyle{empty}
\centering
\vspace*{22mm}

{\Large\color{softgray}Чек-лист\par}
\vspace{8mm}
{\Huge\bfseries\color{darkaccent}Прикладные задачи\par}
\vspace{3mm}
{\Large\color{textgray}Формулы, единицы, квадратные уравнения и отбор корней\par}

\vspace{18mm}
{\large Матвей Горбачев\par}
\vspace{3mm}
{\normalsize 27.09.2026\par}

\vfill
{\large\bfseries\color{darkaccent}Лёвин Артём Александрович\par}
\vspace{2mm}
{\normalsize эксклюзивно для Матвея Горбачева\par}

\vspace{12mm}
\begin{tikzpicture}[remember picture,overlay]
  \draw[accent,line width=1.2pt]
    ($(current page.south west)+(20mm,18mm)$)
    .. controls ($(current page.south west)+(70mm,31mm)$)
    and ($(current page.south east)+(-70mm,7mm)$)
    .. ($(current page.south east)+(-20mm,18mm)$);
\end{tikzpicture}
\end{titlepage}

\section*{1. Главная идея}
\sectionline

Прикладная задача обычно выглядит длиннее, чем её математическое ядро. Ваша цель --- быстро перевести текст в последовательность
\[
\text{данные}
\;\longrightarrow\;
\text{формула}
\;\longrightarrow\;
\text{уравнение или неравенство}
\;\longrightarrow\;
\text{проверка смысла ответа}.
\]

\key{Главный навык:} не считать сразу. Сначала определить, что обозначает каждая величина, какие данные даны явно, какие приходится извлечь из текста, и что именно требуется найти.

\section*{2. Что обязательно проверять до вычислений}
\sectionline

\begin{itemize}
  \item \key{Что ищем?} Выпишите искомую величину и её смысл.
  \item \key{Какая формула дана?} Не путайте запись вида $L(T)$ с произведением $L\cdot T$: это значение величины $L$, зависящее от параметра $T$.
  \item \key{Единицы измерения.} Если в одной формуле встречаются метры и миллиметры, секунды и минуты и т.~п., приведите всё к согласованным единицам.
  \item \key{Неявные данные.} Фразы «частота второго гудка больше», «уровень поднялся», «прибор отключается при достижении температуры» содержат математическую информацию, которую нужно отдельно записать.
  \item \key{Допустимость.} В прикладной задаче математический корень может не иметь физического или смыслового значения.
\end{itemize}

\section*{3. Технические приёмы, которые экономят вычисления}
\sectionline

\subsection*{3.1. Убирайте неудобные десятичные дроби}

Например,
\[
1{,}2\cdot10^{-5}=12\cdot10^{-6},
\qquad
0{,}625=625\cdot10^{-3}.
\]

Так легче сокращать множители и степени десятки, не теряя точность.

\subsection*{3.2. Избавляйтесь от дробей в уравнении}

Если все знаменатели числовые, домножьте обе части на их общий знаменатель. Например,
\[
10=-\frac{x^2}{60}+\frac{7x}{6}
\]
удобно умножить на $60$ и получить
\[
x^2-70x+600=0.
\]

\subsection*{3.3. Рисунок полезен, когда текст описывает процесс}

Если говорится о падении, подъёме уровня воды, движении мяча или расстояниях до линзы, короткая схема часто сразу показывает, какую разность или какой интервал нужно искать.

\section*{4. Типовые модели занятия}
\sectionline

\subsection*{4.1. Изменение величины по формуле}

Если
\[
L=L_0(1+\alpha T),
\]
то при известном удлинении $\Delta L=L-L_0$ удобно сначала записать
\[
L=L_0+\Delta L,
\]
а затем подставить в исходную формулу.

\textbf{Пример.}
\[
L_0=10\text{ м},
\qquad
\alpha=1{,}2\cdot10^{-5},
\qquad
\Delta L=3\text{ мм}=3\cdot10^{-3}\text{ м}.
\]

Тогда
\[
10+3\cdot10^{-3}
=
10\left(1+12\cdot10^{-6}T\right),
\]
откуда
\[
T=25^{\circ}\text{C}.
\]

\subsection*{4.2. Разность двух состояний}

Если путь свободного падения задаётся формулой
\[
h=5t^2,
\]
а время полёта уменьшилось с $t_1$ до $t_2$, то изменение уровня
\[
\Delta h=5t_1^2-5t_2^2.
\]

При $t_1=0{,}6$ с и $t_2=0{,}4$ с:
\[
\Delta h
=
5(0{,}36-0{,}16)
=
1\text{ м}.
\]

\subsection*{4.3. Нулевая величина может упростить формулу}

Если
\[
P=m\left(\frac{v^2}{L}-g\right)
\]
и по условию $P=0$, а $m\ne0$, то
\[
\frac{v^2}{L}-g=0,
\qquad
v=\sqrt{Lg}.
\]

При $L=0{,}625$ м и $g=10$ м/с$^2$:
\[
v^2=6{,}25.
\]

Так как $v$ --- скорость, берём неотрицательное значение:
\[
v=2{,}5\text{ м/с}.
\]

\subsection*{4.4. Два корня --- два момента или два положения}

Квадратное уравнение часто описывает два момента времени или два положения объекта. Перед записью ответа спросите:

\begin{itemize}
  \item нужен первый или второй момент;
  \item нужна разность моментов;
  \item нужен наибольший или наименьший корень;
  \item не завершился ли процесс раньше второго корня.
\end{itemize}

Например, если прибор отключается при первом достижении предельной температуры, из двух положительных корней выбирают \textbf{первый по времени}.

\subsection*{4.5. «Не менее» и «не более» задают неравенство}

Равенство удобно использовать для поиска \textbf{границ}, но окончательный ответ определяется исходным знаком неравенства.

Если высота мяча $h(t)$ должна удовлетворять условию
\[
h(t)\ge h_0,
\]
сначала найдите моменты пересечения из
\[
h(t)=h_0.
\]

Затем определите, где график находится на уровне $h_0$ или выше. Для параболы это можно сделать по направлению ветвей, методом интервалов или проверочной точкой.

Записывать в ответ только корни уравнения здесь недостаточно: вопрос обычно требует \textbf{промежуток времени} или его длину.

\section*{5. Отбор корней по смыслу}
\sectionline

После решения уравнения каждый корень проверяйте по трём фильтрам:

\begin{enumerate}
  \item \key{математический:} удовлетворяет ли исходному уравнению и ограничениям;
  \item \key{предметный:} имеет ли смысл отрицательное время, отрицательная длина, значение после завершения процесса;
  \item \key{формулировочный:} спрашивают первый момент, последний момент, продолжительность, минимум, максимум или конкретное положение.
\end{enumerate}

\section*{6. Задачи на зависимость величин}
\sectionline

Иногда ответ нельзя получить простой подстановкой крайнего значения. Нужно понять, \textbf{как одна величина зависит от другой}.

Например,
\[
\frac1{d_1}+\frac1{d_2}=\frac1f
\]
даёт
\[
\frac1{d_1}=\frac1f-\frac1{d_2}.
\]

Если $f$ фиксировано и нужно минимизировать $d_1$, то дробь $1/d_1$ должна стать максимальной.

При увеличении положительного $d_2$ величина $1/d_2$ уменьшается, поэтому правая часть возрастает, а $d_1$ уменьшается. Значит, для минимума $d_1$ берут \textbf{наибольшее допустимое $d_2$}.

\textbf{Пример.} Пусть
\[
f=30\text{ см},
\qquad
d_2\in[180;210]\text{ см},
\qquad
d_1\in[30;50]\text{ см}.
\]

Для минимума $d_1$ берём $d_2=210$ см:
\[
\frac1{d_1}
=
\frac1{30}-\frac1{210}
=
\frac6{210}
=
\frac1{35},
\qquad
d_1=35\text{ см}.
\]

Полученное значение попадает в заданный диапазон $[30;50]$ см.

\section*{7. Проценты и сокращение неизвестной величины}
\sectionline

Если одна величина составляет $p\%$ от другой,
\[
A=\frac{p}{100}B.
\]

В формулах вида
\[
I=\frac{\mathcal E}{R+r},
\qquad
I_{\text{кз}}=\frac{\mathcal E}{r}
\]
общий ненулевой множитель $\mathcal E$ можно сократить после корректной подстановки отношения $I$ и $I_{\text{кз}}$.

\textbf{Пример.} При $r=1$ Ом и
\[
I=20\%\,I_{\text{кз}}
\]
получаем
\[
\frac{\mathcal E}{R+1}
=
\frac15\cdot\frac{\mathcal E}{1}
\quad\Longrightarrow\quad
\frac1{R+1}
=
\frac15
\quad\Longrightarrow\quad
R=4\text{ Ом}.
\]

\clearpage
\thispagestyle{plain}

\begin{center}
{\Large\bfseries\color{darkaccent}Алгоритм решения прикладной задачи}
\end{center}

\vspace{6mm}

\begin{center}
\begin{tikzpicture}[node distance=7mm]
  \node (start) [startstop] {Прочитайте условие целиком};

  \node (p1) [process,below=of start]
  {Определите искомую величину и расшифруйте обозначения в формуле};

  \node (p2) [process,below=of p1]
  {Извлеките явные и неявные данные; при необходимости сделайте короткий рисунок};

  \node (p3) [process,below=of p2]
  {Приведите единицы к согласованному виду и упростите запись чисел};

  \node (p4) [process,below=of p3]
  {Составьте уравнение или неравенство и выполните вычисления};

  \node (p5) [process,below=of p4]
  {Проверьте все найденные значения: исходное условие, ограничения, единицы и физический смысл};

  \node (d1) [decision,below=10mm of p5]
  {Вопрос требует выбрать момент, минимум, максимум или интервал?};

  \node (p6) [process,below=12mm of d1]
  {Выберите нужное допустимое значение, разность моментов или нужный промежуток};

  \node (finish) [startstop,below=9mm of p6]
  {Сформулируйте ответ с единицами};

  \draw[arrow] (start) -- (p1);
  \draw[arrow] (p1) -- (p2);
  \draw[arrow] (p2) -- (p3);
  \draw[arrow] (p3) -- (p4);
  \draw[arrow] (p4) -- (p5);
  \draw[arrow] (p5) -- (d1);

  \draw[arrow]
    (d1.south)
    --
    node[right,font=\small]{да}
    (p6.north);

  \draw[arrow]
    (d1.east)
    --
    ++(24mm,0)
    |-
    node[pos=0.18,above,font=\small]{нет}
    (finish.east);

  \draw[arrow] (p6) -- (finish);
\end{tikzpicture}
\end{center}

\clearpage

\section*{8. Тренировочный блок --- Путь А}
\sectionline

\textbf{Решите 10 задач.}

Решения оформляйте кратко, но обязательно показывайте переход от текста к формуле и финальную проверку ответа.

\begin{enumerate}

  \item Длина стержня при нагревании описывается формулой
  \[
  L=L_0(1+\alpha T).
  \]
  При $L_0=12$ м и $\alpha=1{,}5\cdot10^{-5}$ стержень удлинился на $5{,}4$ мм. Найдите $T$.

  \item Глубина, пройденная падающим камнем, задаётся формулой
  \[
  h=5t^2.
  \]
  До дождя время падения камня до воды было $0{,}8$ с, после дождя --- $0{,}6$ с. На сколько метров поднялся уровень воды?

  \item Сила давления описывается формулой
  \[
  P=m\left(\frac{v^2}{L}-g\right).
  \]
  Найдите скорость $v$, если
  \[
  P=0,
  \qquad
  L=0{,}9\text{ м},
  \qquad
  g=10\text{ м/с}^2,
  \qquad
  m>0.
  \]

  \item Частоты связаны формулой
  \[
  f=\frac{f_0}{1-v/c}.
  \]
  Первый гудок имеет частоту $f_0=490$ Гц, а частота второго гудка на $10$ Гц больше. Скорость звука $c=300$ м/с. Найдите скорость приближения $v$.

  \item Высота объекта зависит от расстояния $x\ge0$ по формуле
  \[
  y=-\frac{x^2}{48}+\frac54x.
  \]
  При каких значениях $x$ объект находится на высоте $12$ м? В ответ запишите \textbf{наибольшее} значение $x$.

  \item Высота мяча задаётся формулой
  \[
  h(t)=-4t^2+8t+1,
  \qquad
  t\ge0.
  \]
  Сколько секунд мяч находится \textbf{не ниже} отметки $4$ м?

  \item Температура прибора меняется по закону
  \[
  T(t)=1400+200t-10t^2.
  \]
  При первом достижении температуры $1760$ прибор автоматически отключается. Через сколько секунд произойдёт отключение?

  \item Сила тока задаётся формулами
  \[
  I=\frac{\mathcal E}{R+r},
  \qquad
  I_{\text{кз}}=\frac{\mathcal E}{r}.
  \]
  Известно, что $r=2$ Ом и
  \[
  I=25\%\,I_{\text{кз}}.
  \]
  Найдите $R$.

  \item Для линзы выполняется формула
  \[
  \frac1{d_1}+\frac1{d_2}=\frac1{24}.
  \]
  Расстояние $d_2$ может изменяться от $120$ до $168$ см. Найдите наименьшее возможное значение $d_1$.

  \item Высота тела над уровнем старта описывается формулой
  \[
  h(t)=20t-5t^2,
  \qquad
  t\ge0.
  \]
  Найдите промежуток времени, на котором
  \[
  h(t)\ge15\text{ м}.
  \]

\end{enumerate}

\clearpage

\section*{9. Ответы}
\sectionline

\renewcommand{\arraystretch}{1.35}

\begin{center}
\begin{tabularx}{\linewidth}{@{}>{\bfseries}c X@{}}
\toprule
№ & Ответ \\
\midrule
1 & $30^{\circ}\text{C}$ \\
2 & $1{,}4$ м \\
3 & $3$ м/с \\
4 & $6$ м/с \\
5 & $48$ м \\
6 & $1$ с \\
7 & $2$ с \\
8 & $6$ Ом \\
9 & $28$ см \\
10 & $t\in[1;3]$ с \\
\bottomrule
\end{tabularx}
\end{center}

\vfill

{\small\color{softgray}
Перед сдачей каждого ответа проверьте: единицы, ограничения, смысл найденных значений и точное соответствие вопросу задачи.
}

\end{document}